\documentclass[11pt,a4paper]{exam}
\noprintanswers
%\printanswers

\usepackage[T1]{fontenc}
\usepackage{lmodern}
\usepackage{amsmath,amssymb}
\usepackage[margin=2cm,includeheadfoot,headheight=15pt,headsep=18pt]{geometry}
\usepackage{needspace}
\usepackage[hidelinks]{hyperref}

\newcommand{\CourseTitle}{Topics in Math. Sci. VII 2026}
\newcommand{\AssignmentNumber}{3}
\newcommand{\Deadline}{to be announced}

\pagestyle{headandfoot}
\header{}{\small\CourseTitle}
       {\small Assignment \AssignmentNumber}
\headrule
\footer{}{\small Submit via email.  Deadline: \Deadline}{\small\thepage}

\newcommand{\kk}{\Bbbk}
\newcommand{\bbD}{\mathbb D}
\newcommand{\grstar}{(*)}
\DeclareMathOperator{\Hom}{Hom}
\DeclareMathOperator{\End}{End}
\DeclareMathOperator{\rad}{rad}
\newcommand{\SymAlg}{\mathsf{S}}
\newcommand{\op}{\mathrm{op}}
\newcommand{\id}{\mathrm{id}}
\DeclareMathOperator{\Ext}{Ext}
\DeclareMathOperator{\ext}{ext}
\DeclareMathOperator{\Ker}{Ker}
\let\Im\relax
\DeclareMathOperator{\Im}{Im}
\DeclareMathOperator{\grtop}{grtop}
\newcommand{\Z}{\mathbb Z}

\qformat{\textbf{Exercise \thequestion.\ \thequestiontitle}\hfill}
\renewcommand{\questionshook}{%
  \setlength{\leftmargin}{0pt}%
  \setlength{\labelwidth}{-\labelsep}%
  \setlength{\itemsep}{13pt}%
}
\renewcommand{\thepartno}{\arabic{partno}}
\renewcommand{\partlabel}{(\thepartno)}
\renewcommand{\partshook}{\setlength{\itemsep}{4pt}}
\setlength{\parindent}{0pt}
\setlength{\parskip}{4pt}
\setlength{\emergencystretch}{1em}
\raggedbottom

\newcommand{\comm}[1]{}


% DRAFT: collects exercises moved out of Assignment 2; to be rebalanced later.
\begin{document}

Complete at least 2 exercises.

Throughout, $\kk$ is a field and all modules are graded right modules.
A graded algebra $A=\bigoplus_{i\geq0}A_i$ is locally finite-dimensional with $S=A_0$ semisimple; in all explicit examples $S=\bigoplus_i\kk e_i$ with coordinate idempotents $e_i$.

\begin{questions}

\titledquestion{}
\begin{parts}
\part Let $S=\bigoplus_{i=1}^3\kk e_i$ and let $V=\kk a_1\oplus\kk a_2\oplus\kk b_1\oplus\kk b_2$ with
\[
  a_1\in e_1Ve_2,\quad a_2\in e_2Ve_3,\quad b_1\in e_2Ve_1,\quad b_2\in e_3Ve_2,
\]
and let $A=T_S(V)/(R)$ with $R=\operatorname{span}\{a_1b_1,\ b_2a_2,\ b_1a_1-a_2b_2\}$.
\begin{subparts}
\subpart Check that
\[
  h_A(t)=\begin{pmatrix}1&t&t^2\\t&1+t^2&t\\t^2&t&1\end{pmatrix}.
\]
\begin{solution}
The composable words of length $2$ are $a_1b_1,\ a_1a_2,\ b_1a_1,\ a_2b_2,\ b_2b_1,\ b_2a_2$ (six), and the three relations are linearly independent, so $\dim A_2=3$, spanned by $a_1a_2\in e_1Ae_3$, $b_2b_1\in e_3Ae_1$ and $b_1a_1=a_2b_2\in e_2Ae_2$.
Degree $1$ gives the arrows, and degree $0$ gives $I$.
\end{solution}
\subpart Show that $A_d=0$ for $d\geq3$.
\begin{solution}
In degree $3$ every composable word contains $a_1b_1$, $b_2a_2$ or (after substituting $b_1a_1=a_2b_2$) one of them: for example $a_1a_2b_2=a_1b_1a_1=0$, $b_1a_1a_2=a_2b_2a_2=0$, $b_2b_1a_1=b_2a_2b_2=0$, $a_2b_2b_1=b_1a_1b_1=0$; the remaining words are treated the same way.
So $A_3=0$ and $A_d=0$ for $d\geq3$.
\end{solution}
\subpart Using (a), (b) and Assignment~2, Exercise~3, part~(2), prove that $A$ is not Koszul.
\begin{solution}
Write $h_A(t)=I+tC+t^2D$ with $C=\begin{pmatrix}0&1&0\\1&0&1\\0&1&0\end{pmatrix}$ and $D=\begin{pmatrix}0&0&1\\0&1&0\\1&0&0\end{pmatrix}$.
One checks $C^2=I+D$, $DC=CD=C$.
The inverse $\sum_ku_kt^k$ satisfies $u_0=I$, $u_k=-Cu_{k-1}-Du_{k-2}$, so
\[
u_1=-C,\qquad u_2=C^2-D=I,\qquad u_3=-C+DC=0,\qquad u_4=-D .
\]
By that part, if $A$ were Koszul then $u_4$ would have non-negative entries, but $u_4=-D$ does not.
\end{solution}
\end{subparts}

\part Let $S=\bigoplus_{x=1}^n\kk e_x$, let $V$ be a finite-dimensional $S$-bimodule in degree $1$, and let $A=T_S(V)$.
Prove that the multiplication map $V\otimes_SA\to A_+$ is an isomorphism of right $A$-modules, and deduce that $S$ has the minimal graded projective resolution
\[
0\to(V\otimes_SA)(-1)\to A\to S\to0,
\]
so that $A$ is Koszul.
\begin{solution}
The tensor algebra has $A_i=V^{\otimes_Si}$, so the multiplication map $V\otimes_SA_{i-1}\to A_i$ is the identity of $V^{\otimes_Si}$ for $i\geq1$.
Hence $V\otimes_SA\to A_+$, $v\otimes a\mapsto va$, is bijective; it is right $A$-linear and raises degrees by $1$, so $A_+\cong(V\otimes_SA)(-1)$.
This module is projective: every $S$-module is projective, so $V\otimes_SA$ is a direct summand of a free $A$-module.
Thus $0\to A_+\to A\to S\to0$ is a graded projective resolution, with $A_+\subseteq AA_+$, so it is minimal.
Its terms are generated in degrees $0$ and $1$, so it is linear and $A$ is Koszul.
\end{solution}

\part Let $S=\bigoplus_{i=1}^n\kk e_i$, let $V$ be a finite-dimensional $S$-bimodule in degree $1$, put $C_{ij}=\dim_\kk e_iVe_j$, and let $A=T_S(V)/(V\otimes_SV)=S\oplus V$.
\begin{subparts}
\subpart Prove that $A^!=T_S(D(V))$.
\begin{solution}
The relation space is $R=V\otimes_SV$; since $D(V)\otimes_SD(V)\cong D(V\otimes_SV)$, the annihilator of $R$ is $R^\perp=0$ and $A^!=T_S(D(V))$.
\end{solution}
\subpart Prove that $h_A(t)=I+tC$ and $h_{A^!}(t)=(I-tC^{\mathsf T})^{-1}$.
\begin{solution}
$\dim e_iA_0e_j=\delta_{ij}$ and $\dim e_iA_1e_j=C_{ij}$ give $h_A=I+tC$.
For $D(V)$, $\dim e_iD(V)e_j=\dim e_jVe_i$, so its dimension matrix is $C^{\mathsf T}$, and dimension matrices multiply under $\otimes_S$, so $h_{A^!}=\sum_ht^h(C^{\mathsf T})^h=(I-tC^{\mathsf T})^{-1}$.
\end{solution}
\subpart Deduce that $h_A(t)\,h_{A^!}(-t)^{\mathsf T}=I$.
\begin{solution}
Finally $h_{A^!}(-t)^{\mathsf T}=(I+tC)^{-1}=h_A(t)^{-1}$.
\end{solution}
\subpart Prove that the minimal graded projective resolution of $S$ has $P_h\cong(V^{\otimes_Sh}\otimes_SA)(-h)$ for all $h\geq0$.

\emph{Hint.} Show that $\Ker d_0=V$ and that $VA_+=0$, so each kernel is killed by $A_+$.
\begin{solution}
$\Ker d_0=A_+=V$, and $VA_+=VV=0$, so $V$ is a module with $A_+$ acting by zero, that is, a graded $S$-module in degree $1$.
Put $Y_h=V^{\otimes_Sh}$ in degree $h$, killed by $A_+$.
Then $P_h:=(Y_h\otimes_SA)(-h)\to Y_h$, $k\otimes a\mapsto ka$, is surjective, and as $Y_h\otimes_SA=Y_h\otimes_SS\oplus Y_h\otimes_SV$ with $k\otimes v\mapsto kv=0$, its kernel is $Y_h\otimes_SV=V^{\otimes_S(h+1)}=Y_{h+1}$ (in degree $h+1$), again killed by $A_+$.
In each step $\Ker\subseteq P_hA_+$ since $\Ker=Y_h\otimes V\subseteq P_hA_+$, so this is the minimal graded projective resolution, with $P_h=(V^{\otimes_Sh}\otimes_SA)(-h)$ generated in degree $h$.
\end{solution}
\subpart Deduce that $A$ is Koszul.
\begin{solution}
The resolution in~(d) is linear, since $P_h$ is generated in degree $h$.
\end{solution}
\end{subparts}

\part Let $A=\kk[u,v]/(uv)$ with $\deg u=\deg v=1$.
\begin{subparts}
\subpart Prove that $ua+vb=0$ in $A$ if and only if $a\in vA$ and $b\in uA$.
\begin{solution}
A basis of $A$ is $1,u^i,v^j$ for $i,j\geq1$.
For $a=a_0+a_u+a_v$ with $a_u\in\operatorname{span}\{u^i\}_{i\geq1}$ and $a_v\in\operatorname{span}\{v^j\}_{j\geq1}$ we have $ua=a_0u+ua_u$, and similarly for $vb$.
So $ua+vb=0$ forces $a_0u+ua_u=0$ and $b_0v+vb_v=0$ (the $u$-powers and $v$-powers are independent), hence $a_0=a_u=0$ and $b_0=b_v=0$, that is, $a\in vA$ and $b\in uA$.
The converse is clear because $uv=0$.
\end{solution}
\subpart Show that the minimal graded projective resolution of $\kk=A/A_+$ has $P_0=A$ and $P_h\cong A(-h)^{\oplus2}$ for all $h\geq1$.
\begin{solution}
Likewise $va=0$ if and only if $a\in uA$, and $ua=0$ if and only if $a\in vA$.
Now $P_0=A$ and $\Ker d_0=A_+=uA+vA$ is generated by $u,v$, so $P_1=A(-1)^{\oplus2}$ with $d_1(a,b)=ua+vb$ is a surjection onto $\Ker d_0$.
Its kernel is $\{(va',ub')\}$, generated in degree $2$ by $(v,0)$ and $(0,u)$, so $P_2=A(-2)^{\oplus2}$ with $d_2(a',b')=(va',ub')$.
Then $\Ker d_2=\{(a',b'):va'=0,ub'=0\}=uA\oplus vA$, generated in degree $3$ by $(u,0)$ and $(0,v)$, so $P_3=A(-3)^{\oplus2}$ with $d_3(a'',b'')=(ua'',vb'')$.
Finally $\Ker d_3=vA\oplus uA$ has the same form as $\Ker d_1$, so the pattern repeats with period $2$.
Thus $P_h\cong A(-h)^{\oplus2}$ for all $h\geq1$, and all differentials map into $PA_+$, so the resolution is minimal.
\end{solution}
\subpart Deduce that $A$ is Koszul.
\begin{solution}
The resolution in~(b) is linear, since $P_h$ is generated in degree $h$, so $A$ is Koszul.
\end{solution}
\end{subparts}

\end{parts}

\end{questions}

\end{document}
